\documentclass[11pt]{article}
\usepackage[utf8]{inputenc}
\usepackage[T1]{fontenc}
\usepackage[ngerman]{babel}
\usepackage[margin=2.5cm]{geometry}

\title{Algorithm Engineering\\Woche 2, Aufgabe 3}
\author{Jan Vogt}
\date{}

\begin{document}

\maketitle

\section*{Coding Warmup (Folie 20)}

Alle Programme liegen im selben Ordner und wurden auf meinem Laptop (Apple M3) mit
\texttt{g++-16 -fopenmp} kompiliert und ausgeführt.

\subsection*{1. Race Condition mit std::mutex beheben}

Auf Folie 13 schreiben alle Threads am Ende gleichzeitig in die gemeinsame Variable \texttt{pi}.
\texttt{pi += ...} besteht aber aus Lesen, Addieren und Schreiben. Wenn zwei Threads gleichzeitig
lesen, geht der Beitrag von einem der beiden verloren. In \texttt{race\_condition.cpp} wird das
Update deshalb mit einem \texttt{std::mutex} geschützt:
\begin{verbatim}
lock_guard<mutex> lock(pi_mutex);
pi += sum_local * 4 * width;
\end{verbatim}
Mit 100.000.000 Schritten ist der Fehler auch ohne Mutex in 20 Läufen nie aufgetreten, weil die
Threads zu unterschiedlichen Zeiten fertig werden. Deshalb habe ich beide Versionen mit nur 1000
Schritten jeweils 2000 mal laufen lassen:

\begin{center}
\begin{tabular}{lrl}
Version & Anzahl & Ergebnis \\
\hline
ohne Mutex & 1960 & 3.1415927369231262 \\
 & 3 & 3.1415927369231271 \\
 & 18 & 2.3564443034735985 \\
 & 18 & 2.3559443019110908 \\
 & 1 & 2.3554448015985905 \\
\hline
mit Mutex & 1994 & 3.1415927369231262 \\
 & 6 & 3.1415927369231271 \\
\end{tabular}
\end{center}

Ohne Mutex war das Ergebnis in 37 Läufen falsch. Die falschen Werte liegen alle bei ungefähr
$\frac{3}{4}\pi$, es ist also genau der Beitrag von einem Thread verloren gegangen. Mit Mutex war
das Ergebnis immer richtig. Dass sich in manchen Läufen die letzte Stelle unterscheidet, liegt daran,
dass die Threads ihre Teilsummen in unterschiedlicher Reihenfolge addieren.

\subsection*{2. Schedule für das Denoising-Programm}

Auf Folie 16 braucht Bild \texttt{i} eine Arbeit proportional zu $i^2$. Ohne Schedule verteilt
OpenMP die Schleife statisch in vier gleich große Blöcke. Weil die Schleife rückwärts läuft,
bekommt Thread 0 die Bilder 63 bis 48 und hat damit viel mehr Arbeit als die anderen. Gemessen
habe ich mit 4 Threads (jeweils 5 Läufe):

\begin{center}
\begin{tabular}{lr}
Schedule & Zeit [s] \\
\hline
ohne Schedule & 0.503 bis 0.507 \\
\texttt{static} & 0.504 bis 0.506 \\
\texttt{static, 1} & 0.240 bis 0.320 \\
\texttt{dynamic} & 0.224 bis 0.237 \\
\texttt{dynamic, 1} & 0.224 bis 0.230 \\
\texttt{dynamic, 4} & 0.225 bis 0.239 \\
\texttt{guided} & 0.505 bis 0.513 \\
\end{tabular}
\end{center}

Mit \texttt{schedule(dynamic, 1)} holt sich jeder Thread immer das nächste Bild, sobald er fertig
ist. Dadurch ist das Programm etwa 2,2 mal schneller. \texttt{guided} bringt hier nichts, weil es
am Anfang große Blöcke verteilt und der erste Block genau die schwersten Bilder enthält.

\subsection*{3. Welcher Schedule erzeugt das Muster?}

Das Muster auf Folie 20 entsteht mit \texttt{schedule(static, 22)}. Die 32 Iterationen werden in
Blöcke mit 22 Iterationen aufgeteilt. Thread 0 bekommt den ersten Block (0 bis 21), Thread 1 den
Rest (22 bis 31) und Thread 2 und 3 bekommen gar nichts. Das habe ich mit
\texttt{schedule\_clause.cpp} getestet und es kam in jedem Lauf genau diese Aufteilung heraus. Mit
\texttt{dynamic, 22} entstehen zwar dieselben zwei Blöcke, aber welcher Thread sie bekommt, ändert
sich bei jedem Lauf.

\end{document}
